\chapter{Category theory}\label{ch:category_theory}

Category theory uses relationships between objects to study the objects themselves. This shifts the focus from an object's intrinsic properties based on set membership and function application.

This shift is evident from the following diagram from \cref{thm:adjunction_from_universal_morphism}, which expresses how the constructions in the proof relate to each other:
\begin{equation*}
  \begin{aligned}
    \includegraphics[page=7]{output/thm__category_completeness_via_products_and_equalizers}
  \end{aligned}
\end{equation*}

We will define categories in \fullref{sec:categories} as \hyperref[def:directed_multigraph]{directed multigraphs} with additional structure. Constructing the essential categories involves some size considerations, which we solve by assuming the existence of a \hyperref[def:grothendieck_universe]{Grothendieck universe} (i.e. axiom \logic{U} in \hyperref[def:zfcu]{\logic{ZFC+U}}).

The most essential topic beyond categories themselves are functors, which describe relationships between categories. These are discussed in \fullref{sec:functors}. Natural transformations between functors are discussed in \fullref{sec:natural_transformations} and \fullref{sec:yoneda_embeddings}.

Beyond that, \fullref{sec:adjoint_functors} and \fullref{sec:categorical_limits} contain some topics important for us, and \fullref{sec:category_similarity} fills some gaps.
